\section{詳細解説}

\subsection{脚の質量と AMBAC の強さ}

総質量 76 kg・腕 4 kg を固定し、脚の質量を全質量の 2〜40\% で変えて、1 周あたりの回転の上界を求めた（図~\ref{fig:D12_leg_mass.png}）。
脚が軽い（2\%）と 2.2°、人間型（16\%）で 9.6°、重い（40\%）と 14.3° になる。
上界は脚の質量とともに単調に増えるが、増え方は次第に緩やかになる（脚が重くなると、胴体の側が相対的に軽くなり、$\Gamma$ が 1 に近づいて頭打ちになるため）。
\textbf{宇宙で使う機体なら、推進剤タンクや装甲などの重い装備を脚に置くほど、AMBAC が強くなる}。

\DFIG{D12_leg_mass.png}{宇宙で、脚の質量の割合と 1 周あたりの回転の上界。破線は人間型（脚 12 kg）。}

\subsection{腰の戦略の最適化の定式化}

二本脚の荷物運びでは、腕の軌道（関節ごとに 4 つの補正係数）に加えて、上体の前後傾 $\lambda(\tau)$ を最適化した。
\begin{equation}
\lambda(\tau) = \lambda_0 + (\lambda_1-\lambda_0)\,m(\tau) + \sum_{k=0}^{3} d_k\sin^2(\pi\tau)\cos(k\pi\tau),
\end{equation}
ここで $m(\tau)$ は最小躍度のプロファイル、$\lambda_0$, $\lambda_1$ は動作の開始・終了時の静的な傾きである。
最初は $\lambda_0=\lambda_1=0$ に固定していたため、静止した状態で後ろに反る姿勢が取れず、解が見つからなかった（付録~\ref{app:num} の修正の記録）。
制約は、ZMP（足裏の内側 2 cm 以内）、滑り、腕の関節の速度・加速度、上体の傾き（$\pm40^\circ$）・速度（3 rad/s）・加速度（15 rad/s$^2$）、
床から浮かないこと（垂直反力 $\ge$ 自重の 30\%）である。

\begin{example}[なぜ静止していても転倒するか]
二本脚で、荷物 10 kg を前方 0.64 m に持って静止すると、Sim06 と同じ計算で重心の水平位置（＝静的な ZMP）は約 0.13 m になる。
足裏の支持範囲は $\pm0.125$ m なので、安全余裕を考えなくても範囲を超える。
上体（56 kg）を後ろに $\lambda$ だけ傾けると、上体の重心（股関節から 0.25 m）が $0.25\sin\lambda$ だけ後ろに動き、全体の重心は
$56\times0.25\sin\lambda/97$ だけ戻る。$\lambda=-40^\circ$ では約 0.09 m 戻り、ZMP を足裏の中に収められる。
\end{example}

\subsection{関節の数とコスト}

\begin{longtable}{p{0.30\linewidth}p{0.30\linewidth}p{0.30\linewidth}}
\toprule
& 戦車型＋腰 & 二本脚（人間型） \\ \midrule \endhead
移動に使う関節 & 履帯 2 & 片脚 6 × 2 = 12 \\
姿勢に使う関節 & 腰 1〜2 & 腰・体幹 \\
立っているだけの制御 & 不要 & 常に必要 \\
本報告書での平地の荷物運び & 0.64 s & 0.69 s（腰の戦略が必須） \\
本報告書での段差（人間大） & 0.25 m & 0.60 m \\
\bottomrule
\end{longtable}

\section*{考察課題}
\begin{exercise}
二本脚が足を前後に開いて立つ（支持範囲を前後に広げる）と、荷物運びの結果はどう変わるか。腰の戦略は不要になるか。
\end{exercise}
\begin{exercise}
宇宙で脚を重くすると AMBAC は強くなるが、関節を動かすのに必要なトルクも増える（第~\ref{ch:dyn}章）。脚の質量の最適な割合を決めるには、何を評価すればよいか。
\end{exercise}
